%!TEX root = ../main.tex
%!TEX program = xelatex
% Chapter 6: Words 1251 - 1500
\chapter{Complex Descriptions}
\label{chap:chapter6}

{\small\sffamily\color{mutedtext} Descriptive adjectives, compound expressions, specialized terminology, and literary phrasing.\par}

\vspace{0.4em}
\markboth{Chapter 6: Complex Descriptions}{Words 1251--1500}

\begin{multicols}{2}
\vocentry{1251}{Вслед}{vsled}{After, behind, following}{Она с грустью смотрела вслед уходящему поезду.}{She was looking after the leaving train with sadness.}
\vocentry{1252}{Сотрудник}{sat'rudnik}{An employee}{Каждый сотрудник компании получил подарок к Рождеству.}{Every employee of the company got a present for Christmas.}
\vocentry{1253}{Внутри}{vnut'ri}{Inside}{Не открывай коробку! Ты не знаешь, что внутри.}{Don’t open the box! You don’t know what’s inside.}
\vocentry{1254}{Собственно}{'sobstvennə}{Actually, as a matter of fact, in fact}{Я, собственно, не собираюсь оставаться здесь навсегда.}{Actually, I’m not going to stay here forever.}
\vocentry{1255}{Родина}{'rodina}{Motherland}{Моя родина находится далеко отсюда, и я очень по ней скучаю.}{My motherland is far away from here and I miss it a lot.}
\vocentry{1256}{Царь}{tsar’}{Tsar}{Николай II – это последний русский царь.}{Nicolas II is the last Russian tsar.}
\vocentry{1257}{Насчёт}{nas'tchjot}{About, concerning (colloquial)}{Как насчёт сходить завтра в кафе?}{What about going to the cafe tomorrow?}
\vocentry{1258}{Страница}{stra'nitsa}{A page}{Похоже, это страница из твоего дневника.}{Looks like it’s a page from your diary.}
\vocentry{1259}{Скорость}{'skorast’}{Speed}{Какова максимальная скорость этого автомобиля?}{What’s the maximum speed of this car?}
\vocentry{1260}{Зверь}{zver’}{An animal, a beast; a brute (figurative)}{Этот белый медведь выглядит милым, но это дикий и опасный зверь.}{This polar bear looks cute but it’s a wild and dangerous animal.}
\vocentry{1261}{Помочь}{pa'motch}{To help (perfective)}{Я смогу помочь только после обеда. Пойдёт?}{I’ll be able to help only in the afternoon. Will it do?}
\vocentry{1262}{Заявить}{zaja'vit’}{To declare, to state, to claim}{Он был достаточно смелым, чтобы открыто заявить о своём недовольстве властями.}{He was bold enough to openly declare his dissatisfaction with the authorities.}
\vocentry{1263}{Нечто}{'netchtə}{Something; awesome (colloquial)}{Нечто странное и большое шевелилось в темноте.}{Something big and strange was moving in the darkness.}
\vocentry{1264}{Размер}{raz'mer}{A size}{Тебе не кажется, что мне нужен размер побольше?}{Doesn’t it seem to you that I need a bigger size?}
\vocentry{1265}{Житель}{'zhitel’}{An inhabitant, a resident}{Этот пожилой человек – местный житель острова.}{This elderly man is a local inhabitant of the island.}
\vocentry{1266}{Секретарь}{sekre'tar’}{A secretary}{Мне нужен опытный секретарь со знанием иностранных языков.}{I need an experienced secretary with the knowledge of foreign languages.}
\vocentry{1267}{Городской}{garats'koj}{Urban; town, city (adj)}{Городской темп жизни не для меня.}{The urban pace of life is not for me.}
\vocentry{1268}{Прекрасно}{prek'rasnə}{Beautifully, nice, great}{Ваша дочь прекрасно поёт. Кем она хочет стать в будущем?}{Your daughter sings beautifully. What does she want to be in the future?}
\vocentry{1269}{Торчать}{tar'chat’}{To stick out; to hang around (colloquial)}{Старайся выглядеть опрятно. Рубашка не должна вот так вот торчать из брюк.}{Try to look neat. The shirt shouldn’t stick out of your trousers like this.}
\vocentry{1270}{Многое}{'mnogaje}{Many things, much}{Я многое понял после того, как стал родителем.}{I realized many things after I became a parent.}
\vocentry{1271}{Америка}{a'merica}{America}{Америка была названа в честь Америго Веспуччи.}{America was named in honor of Amerigo Vespucci.}
\vocentry{1272}{Сигарета}{siga'reta}{A cigarette}{У тебя осталась только одна сигарета. Может, пора бросать курить?}{You’ve got only one cigarette left. Maybe it's time to give up smoking?}
\vocentry{1273}{Полностью}{'polnastju}{Completely, entirely, fully}{Я полностью согласна с твоим мнением.}{I completely agree with your opinion.}
\vocentry{1274}{Пиво}{'pivə}{Beer}{Я не люблю пиво. Оно для меня слишком горькое на вкус.}{I don’t like beer. It tastes too bitter to me.}
\vocentry{1275}{Приём}{pri'jom}{A reception; an appointment (very often at the doctor’s)}{Это был роскошный приём! Я под впечатлением!}{That was a luxurious reception! I’m impressed!}
\vocentry{1276}{Противник}{pra'tivnik}{An opponent; an enemy}{Наш противник очень опытный, но я думаю, что мы выиграем завтрашний матч.}{Our opponent is very experienced, but I think we’ll win tomorrow’s match.}
\vocentry{1277}{Образование}{abrazo'vanije}{Education}{Сегодня мы считаем образование необходимостью, а не роскошью.}{Today we consider education to be a necessity, not a luxury.}
\vocentry{1278}{Продукт}{pra'duct}{A product (commodity for sale); a food stuff}{Сейчас наша команда тестирует новый продукт.}{Our team is testing a new product now.}
\vocentry{1279}{Поймать}{paj'mat’}{To catch (both literally and figuratively)}{Я надеюсь поймать такси и вернуться домой до полуночи.}{I hope to catch a taxi and to return home before midnight.}
\vocentry{1280}{Пуля}{'pulja}{A bullet}{К счастью, это была не настоящая пуля, а резиновая.}{Luckily, that wasn’t a real bullet but a rubber one.}
\vocentry{1281}{Книжка}{'knizhka}{A book (more colloquial variant of ‘книга’)}{Это не моя книжка. Я взял её в библиотеке.}{It’s not my book. I borrowed it from the library.}
\vocentry{1282}{Дальний}{'dal’nij}{Far, distant, remote}{Его отправили в командировку на Дальний Восток.}{He was sent on a business trip to the Far East.}
\vocentry{1283}{Экономический}{ɛkana'mitcheskij}{Economic, economical}{В следующем году наша страна примет международный экономический форум.}{Next year our country is hosting an international economic forum.}
\vocentry{1284}{Представитель}{predsta'vitel’}{A representative}{Представитель нашей партии не смог присутствовать на собрании.}{The representative of our party was unable to be present at the meeting.}
\vocentry{1285}{Собирать}{sabi'rat’}{To gather, to collect, to harvest}{Это приложение помогает собирать информацию потребителей.}{This application helps to gather consumers’ information.}
\vocentry{1286}{Видимый}{'vidimyj}{Visible}{Нам нужен видимый результат, а не бесконечные обещания.}{We need a visible result and not endless promises.}
\vocentry{1287}{Производство}{praiz'vodstvə}{Production, manufacture}{Страна расширяет производство натуральных продуктов.}{The country is expanding the production of natural foods.}
\vocentry{1288}{Больница}{bal’nitsa}{A hospital}{Местная больница испытывает недостаток врачей.}{The local hospital is experiencing a lack of doctors.}
\vocentry{1289}{Ресторан}{resta'ran}{A restaurant}{Этот ресторан знаменит своим шеф-поваром.}{This restaurant is famous for its chef.}
\vocentry{1290}{Находить}{naha'dit’}{To find; to consider, to think}{Студенты должны уметь находить информацию самостоятельно.}{Students must be able to find information independently.}
\vocentry{1291}{Диван}{di'van}{A sofa, a couch}{Садись на диван. Он мягче, чем стулья.}{Sit down on the sofa. It’s softer than the chairs.}
\vocentry{1292}{Дыхание}{dy'hanije}{Breathing; breath}{Тяжёлое дыхание пациента волновало врача.}{The patient’s heavy breathing troubled the doctor.}
\vocentry{1293}{Тюрьма}{tjur'ma}{A prison}{Это тюрьма для несовершеннолетних преступников.}{This prison is for underage criminals.}
\vocentry{1294}{Лёгкое}{'ljohkəje}{A lung}{Врачам пришлось удалить ей одно лёгкое, чтобы спасти ей жизнь.}{The doctors had to remove one lung of hers to save her life.}
\vocentry{1295}{Мясо}{'m’asə}{Meat}{Я вегетарианец. Я не ем мясо уже более десяти лет.}{I’m a vegetarian. I haven’t been eating any meat for more than ten years already.}
\vocentry{1296}{Внизу}{vni'zu}{Below, underneath; at the foot of}{Внизу есть список необходимых документов.}{There’s a list of necessary documents below.}
\vocentry{1297}{Кошка}{'koshka}{A cat (female)}{Наша кошка родила четыре милых котёнка.}{Our cat gave birth to four cute kittens.}
\vocentry{1298}{Взрослый}{'vzroslyj}{Grown-up; a grown-up, an adult;}{Ты уже взрослый человек и должен уметь заботиться о себе.}{You are a grown-up person already and must be able to take care of yourself.}
\vocentry{1299}{Отсутствие}{at'sutstvije}{Absence}{Отсутствие чётких инструкций привело к недопониманию.}{The absence of clear instructions led to misunderstanding.}
\vocentry{1300}{Настроение}{nastra'jenije}{Mood}{Сегодня мой день рождения и у меня отличное настроение!}{Today is my birthday and I’m in a great mood!}
\vocentry{1301}{Собрать}{sab'rat’}{To gather, to collect, to harvest (perfective)}{Для опроса нужно собрать и проанализировать много информации.}{It’s needed to gather and analyze lots of information for the survey.}
\vocentry{1302}{Наступить}{nastu'pit’}{To come, to begin; to step on}{Почему лето не может наступить быстрее? Я так устала от холода!}{Why can’t summer come sooner? I’m so tired of the cold!}
\vocentry{1303}{Покой}{pa'koj}{Calm, rest, peace}{После такого стрессового дня на работе мне нужен абсолютный покой.}{After such a stressful day at work I need an absolute calm.}
\vocentry{1304}{Тянуть}{tja'nut’}{To pull, to drag}{Когда клюёт, нужно тянуть удочку сильнее.}{When it bites you should pull the rod harder.}
\vocentry{1305}{Народный}{na'rodnyj}{People’s; national; folk}{Верховный народный суд отказался рассматривать это дело.}{The Supreme People’s Court refused to hear the case.}
\vocentry{1306}{Процент}{pra'tsent}{Percentage, percent}{В африканских странах большой процент населения остаётся неграмотным.}{In African countries a great percentage of the population remains illiterate.}
\vocentry{1307}{Постоянно}{pasta'jannə}{Constantly, always}{Ты постоянно теряешь свои ключи!}{You’re constantly losing your keys!}
\vocentry{1308}{Бросать}{bra'sat’}{To give up; to throw}{Тебе пора бросать есть так много сладкого! Это вредно для здоровья!}{It’s time for you to give up eating so many sweet things! It’s bad for health!}
\vocentry{1309}{Мечтать}{metch'tat’}{To dream}{Я не могла даже мечтать о таком подарке! Огромное спасибо!}{I couldn’t even dream of such a present! Thanks a lot!}
\vocentry{1310}{Захотеть}{zaha'tet’}{To want (to begin to want)}{Тебе нужно только захотеть этого, и у тебя всё получится.}{You should only want it and you’ll succeed.}
\vocentry{1311}{Автобус}{af'tobus}{A bus}{Я опоздала на автобус, теперь мне придётся ждать ещё целый час.}{I’ve missed the bus and now I’ll have to wait for a whole hour.}
\vocentry{1312}{Восток}{vas'tok}{East}{Восток страны был оккупирован врагом.}{The East of the country was occupied by the enemy.}
\vocentry{1313}{Оставлять}{astav'ljat’}{To leave; to abandon}{Я не хочу оставлять её одну в этот трудный период.}{I don’t want to leave her alone during this hard period.}
\vocentry{1314}{Усмехнуться}{usmeh'nutsa}{To grin, to smile}{Его истории всегда заставляют меня усмехнуться.}{His stories always make me grin.}
\vocentry{1315}{Эхо}{'ɛhə}{An echo}{Мы слышали эхо наших голосов с другой стороны долины.}{We heard the echo of our voices from the other side of the valley.}
\vocentry{1316}{Примерно}{pri'mernə}{Approximately, about}{Этот магазин примерно в километре отсюда.}{This shop is approximately a kilometer away from here.}
\vocentry{1317}{Автомобиль}{avtama'bil’}{A car, an automobile}{Это гоночный автомобиль новой модели.}{It’s a new model race car.}
\vocentry{1318}{Внешний}{'vneshnyj}{Outer, external}{Внутренний мир человека важнее, чем внешний.}{The inner world of a person is more important than the outer one.}
\vocentry{1319}{Технический}{teh'nitcheskij}{Technical}{Технический прогресс сильно изменил наш мир.}{Technical progress has greatly changed our world.}
\vocentry{1320}{Пропасть}{pro'past’}{To get lost, to go missing; to disappear}{Группа туристов могла пропасть в горах, но никто не уверен.}{The group of tourists could get lost in the mountains but nobody is sure.}
\vocentry{1321}{Выдержать}{'vyderzhat’}{To stand, to bear, to endure}{Их брак не смог выдержать проверку временем.}{Their marriage was unable to stand the test of time.}
\vocentry{1322}{Миг}{mig}{An instant, a wink, a moment}{Наша жизнь – это лишь миг в истории человечества.}{Our life is just an instant in the history of mankind.}
\vocentry{1323}{Отряд}{at'rjad}{A squad, a unit, a team}{Отряд полиции задержал грабителей прямо на месте преступления.}{The police squad detained the robbers right at the crime scene.}
\vocentry{1324}{Сверху}{'sverhu}{On top; from above; above, over}{Я хочу этот торт с клубникой сверху.}{I want this cake with strawberries on top.}
\vocentry{1325}{Настолько}{nas'tol’kə}{So, so much}{Ребёнок был настолько напуган, что не мог сказать ни слова.}{The child was so scared that he couldn’t say a word.}
\vocentry{1326}{Встречать}{vstre'tchat’}{To meet; to greet, to welcome}{Всегда приятно встречать интересных людей.}{It’s always nice to meet interesting people.}
\vocentry{1327}{Пускать}{pus'kat’}{To let in; to permit}{Они отказались меня пускать.}{They refused to let me in.}
\vocentry{1328}{Ясный}{'jasnyj}{Clear; bright}{Она дала ясный и детальный ответ.}{She gave a clear and detailed answer.}
\vocentry{1329}{Занятие}{za'njatije}{An occupation, a business; a lesson}{Ты бездельничаешь весь день! Найди себе какое-нибудь занятие!}{You’ve been idling all day! Find yourself an occupation!}
\vocentry{1330}{Реальный}{re'al’nyj}{Real}{Это невероятно! Неужели это реальный человек, а не выдумка?}{It’s unbelievable! Can it be a real man and not some fiction?}
\vocentry{1331}{Задний}{'zadnij}{Back, rear}{Наш задний дворик уютный и тихий. Я люблю проводить там время.}{Our back yard is cozy and quiet. I like to spend time there.}
\vocentry{1332}{Европа}{jev'ropa}{Europe}{Европа окружена морями с трёх сторон.}{Europe is surrounded by seas on three sides.}
\vocentry{1333}{Случайно}{slu'tchajnə}{Accidentally, by accident}{Извини, я случайно разбил твою любимую чашку.}{Sorry, I’ve accidentally broken your favorite cup.}
\vocentry{1334}{Артист}{ar'tist}{An artist, a performer}{Народный артист – это почётное звание.}{A people’s artist is an honorary title.}
\vocentry{1335}{Сложный}{'slozhnyj}{Complicated, difficult, complex}{Этот текст слишком сложный для студентов твоего уровня.}{This text is too complicated for the students of your level.}
\vocentry{1336}{Планета}{pla'neta}{A planet}{Юпитер – это самая большая планета Солнечной системы.}{Jupiter is the biggest planet of the solar system.}
\vocentry{1337}{Порог}{pa'rog}{A doorstep, a threshold}{Он отвратительный человек. Я не хочу даже пускать его на свой порог.}{He’s a disgusting person. I don’t even want to let him get to my doorstep.}
\vocentry{1338}{Явление}{jav'lenije}{A phenomenon}{Солнечное затмение – редкое и красивое явление.}{A Solar eclipse is a rare and beautiful phenomenon.}
\vocentry{1339}{Отпустить}{atpus'tit’}{To let go}{Чтобы идти дальше, тебе нужно отпустить своё прошлое.}{To go on you should let go of your past.}
\vocentry{1340}{Красота}{krasa'ta}{Beauty}{Красота этого пейзажа поразительна!}{The beauty of this scenery is stunning!}
\vocentry{1341}{Воскликнуть}{vas'kliknut’}{To exclaim, to cry}{Когда я увидела это платье, мне хотелось воскликнуть: “Вот то, что нужно!”}{When I saw that dress I wanted to exclaim: ‘That’s what I need!’}
\vocentry{1342}{Напротив}{nap'rotiv}{On the contrary}{Я не против твоей идеи, напротив, она мне очень нравится.}{I’m not against your idea; on the contrary, I like it a lot.}
\vocentry{1343}{Дрожать}{dra'zhat’}{To tremble, to shiver, to shake}{Я очень боюсь змей. Только один их вид заставляет меня дрожать.}{I’m very much scared of snakes. The only sight of them makes me tremble.}
\vocentry{1344}{Палатка}{pa'latka}{A tent}{Эта туристическая палатка водонепроницаемая и очень прочная.}{This tourist tent is waterproof and very robust.}
\vocentry{1345}{Выбрать}{'vybrat’}{To choose, to select (perfective)}{Я не смог выбрать правильный ответ на последний вопрос, но всё равно сдал тест.}{I didn’t manage to choose the correct answer to the last question, but still passed the test.}
\vocentry{1346}{Пьяный}{'pjanyj}{Drunk; a drunk man}{Ты пьяный и не можешь садиться за руль.}{You’re drunk and can’t drive.}
\vocentry{1347}{Рада}{'rada}{Rada}{Парламент Украины называется Рада.}{The parliament of Ukraine is called Rada.}
\vocentry{1348}{Доска}{das'ka}{A board, a plate; a blackboard}{Эта разделочная доска сделана из натурального дерева.}{This cutting board is made of natural wood.}
\vocentry{1349}{Пистолет}{pista'let}{A pistol, a handgun}{Этот пистолет – главная улика расследования.}{This pistol is the main piece of evidence of the investigation.}
\vocentry{1350}{Пожать}{pa'zhat’}{To shake; to reap (perfective)}{Позвольте мне пожать Вашу мужественную руку!}{Let me shake your courageous hand.}
\vocentry{1351}{Актёр}{ak'tjor}{An actor}{Этот актёр так талантлив! Каждая его новая роль – это шедевр!}{This actor is so talented! His every new role is a masterpiece!}
\vocentry{1352}{Раздаться}{raz'datsa}{To ring out (often unexpectedly)}{Сейчас должен раздаться телефонный звонок. Я жду его весь день.}{A phone call should ring out now. I’ve been waiting for it the whole day.}
\vocentry{1353}{Зря}{zrja}{In vain, for nothing}{Мы зря пытались изменить его мнение. Он был упрям, как всегда.}{We tried in vain to change his mind. He was stubborn, as usual.}
\vocentry{1354}{Потолок}{pata'lok}{A ceiling}{Потолок в комнате слишком низкий, и поэтому она кажется маленькой.}{The ceiling in the room is too low and that’s why it seems small.}
\vocentry{1355}{Постель}{pas'tel’}{A bed; bedding}{Сегодня я проспала и не успела застелить постель.}{I overslept today and didn’t have time to make my bed.}
\vocentry{1356}{Навстречу}{navst'retchu}{Toward, towards}{Переезд в другой город стал для нас первым шагом навстречу большим переменам.}{Moving to a different city became a first step toward big changes for us.}
\vocentry{1357}{Ближайший}{bli'zhajshij}{Nearest; next}{Скажите пожалуйста, где находится ближайший банк?}{Could you tell me please where the nearest bank is?}
\vocentry{1358}{Сумка}{'sumka}{A bag, a handbag}{Эта сумка компактная и довольно вместительная.}{This bag is compact and quite spacious.}
\vocentry{1359}{Мгновение}{mgna'venije}{An instant, a moment (synonymous to ‘миг’)}{Лето прошло быстро, в одно мгновение.}{The summer passed fast, in an instant.}
\vocentry{1360}{Праздник}{'praznik}{A holiday}{Новый Год – это традиционный семейный праздник.}{New Year is a traditional family holiday.}
\vocentry{1361}{Густо}{'gustə}{Densely, thickly}{Индия – одна из самых густонаселённых стран.}{India is one of the most densely populated countries.}
\vocentry{1362}{Мост}{most}{A bridge}{Мост через реку соединяет две части города.}{The bridge across the river connects the two parts of the city.}
\vocentry{1363}{Энергия}{ɛ'nergija}{Energy, power}{Завтрак для меня – это энергия дня.}{Breakfast is the energy of the day for me.}
\vocentry{1364}{Начальство}{na'tchal’stvə}{Boss, authority (collective)}{Здорово быть фрилансером – ты сам себе начальство.}{It’s great to be a freelancer – you’re your own boss.}
\vocentry{1365}{Сожаление}{sozha'lenije}{Regret, pity}{Я слышу сожаление в твоём голосе. Ты не рада, что приехала сюда?}{I hear regret in your voice. Are you unhappy about coming here?}
\vocentry{1366}{Полтора}{palta'ra}{One and a half}{Я освобожусь через полтора часа.}{I’ll be free in an hour and a half.}
\vocentry{1367}{Объяснять}{abjas'njat’}{To explain, to clarify}{Тебе не нужно объяснять мне, каково это – не иметь работы. Я знаю это по своему опыту.}{You don’t need to explain it to me what it feels like to be out of work. I know it from experience.}
\vocentry{1368}{Обращать}{abrast'chat’}{To pay, to turn}{Не стоит обращать внимание на то, что говорят другие.}{It’s not worth paying attention to what other people say.}
\vocentry{1369}{Сунуть}{'sunut’}{To stick, to poke}{Мои родители используют любую возможность сунуть свой нос в мою личную жизнь.}{My parents make use of any opportunity to stick their noses in my personal life.}
\vocentry{1370}{Пятнадцать}{pjat'nadtsat’}{Fifteen}{Такси приедет через пятнадцать минут. Ты готова?}{The taxi will arrive in fifteen minutes. Are you ready?}
\vocentry{1371}{Разрешить}{razre'shit’}{To allow, to permit, to let}{Я не могу разрешить тебе играть с этим. Этот сувенир мне очень дорог.}{I can’t allow you to play with it. This souvenir is very dear to me.}
\vocentry{1372}{Участок}{u'tchastək}{A stretch, a piece of land}{Этот участок дороги самый опасный.}{This stretch of road is the most dangerous one.}
\vocentry{1373}{Жертва}{'zhertva}{A victim; a sacrifice}{В этой ситуации я всего лишь жертва обстоятельств.}{In this situation I’m just a victim of circumstances.}
\vocentry{1374}{Меч}{mjetch}{A sword}{«Меч короля Артура» – это очень популярная легенда.}{“King Arthur's sword” is a very popular legend.}
\vocentry{1375}{Звонок}{zva'nok}{A call, a ring}{Я жду важный деловой звонок.}{I’m expecting an important business call.}
\vocentry{1376}{Кино}{ki'no}{Cinema; a movie, a film}{Сегодня вечером мы идём в кино. Ты с нами?}{We’re going to the cinema today. Are you with us?}
\vocentry{1377}{Признаться}{priz'natsa}{To admit, to confess}{Я должен признаться, у тебя богатое воображение.}{I must admit you’ve got a rich imagination.}
\vocentry{1378}{Мороз}{ma'roz}{Frost, freezing weather}{Сильный мороз побил посевы.}{The heavy frost damaged the crops.}
\vocentry{1379}{Составлять}{sastav'ljat’}{To constitute, to comprise; to construct, to compose}{Все государственные органы должны составлять единую систему.}{All the state organs must constitute a unified system.}
\vocentry{1380}{Бедный}{'bednyj}{Poor (both financially and defined by other circumstances)}{Мой отец слишком бедный, чтобы позволить себе такую машину.}{My father is too poor to afford such a car.}
\vocentry{1381}{Летний}{'ljetnij}{Summer (adj)}{В августе я еду в летний лагерь.}{I’m going to a summer camp in August.}
\vocentry{1382}{Видать}{vi'dat’}{To see (colloquial)}{Мне отсюда не видать, едет его машина или нет.}{I can’t see from here if his car is coming or not.}
\vocentry{1383}{Адрес}{'adres}{An address}{Пришли мне свой новый адрес в сообщении, пожалуйста.}{Please, send me your new address in a message.}
\vocentry{1384}{Закрытый}{zak'rytyj}{Closed; secret}{Для эксперимента вам понадобится закрытый сосуд.}{You’ll need a closed vessel for the experiment.}
\vocentry{1385}{Горло}{'gorlə}{A throat}{У меня ужасно болит горло, и я едва могу разговаривать.}{I have a terrible sore throat and I can barely talk.}
\vocentry{1386}{Костёр}{kas'tjor}{A fire, a bonfire, a wood fire}{Мы установили палатки, развели костёр и начали готовить.}{We put up our tents, made a fire and started cooking.}
\vocentry{1387}{Жалко}{'zhalkə}{It is a pity, unfortunately}{Жалко, что ты не сможешь прийти ко мне на день рождения.}{It is a pity you won’t be able to come to my birthday party.}
\vocentry{1388}{Кот}{kot}{A cat (male)}{Наш кот не ест ничего, кроме рыбы.}{Our cat doesn’t eat anything except fish.}
\vocentry{1389}{Бык}{byk}{A bull, an ox}{Бык – это упрямое животное, которое может быть довольно опасным.}{A bull is a stubborn animal that can be quite dangerous.}
\vocentry{1390}{Трудный}{'trudnyj}{Hard, difficult}{Ей пришлось сделать трудный выбор между любовью и карьерой.}{She had to make a difficult choice between love and her career.}
\vocentry{1391}{Текст}{tekst}{A text}{Этот простой текст подойдёт начинающим.}{This simple text will be suitable for beginners.}
\vocentry{1392}{Тепло}{tep'lo}{Warmth; it is warm (impersonal about weather)}{Детям больше всего нужно внимание и тепло родителей.}{Children need the attention and warmth of their parents most of all.}
\vocentry{1393}{Занимать}{zani'mat’}{To occupy, to take up; to borrow}{Я думаю, этот шкаф будет занимать слишком много места в нашей спальне.}{I think this wardrobe will occupy too much space in our bedroom.}
\vocentry{1394}{Песок}{pe'sok}{Sand}{Тёплый песок, яркое солнце – день на пляже был идеальным.}{The warm sand, the bright sun – the day at the beach was perfect.}
\vocentry{1395}{Больно}{'bol’nə}{Painful, painfully; bad, badly}{Мне больно это говорить, но я думаю, нам стоит расстаться.}{It’s painful for me to say it but I think we should break up.}
\vocentry{1396}{Сомнение}{sam'nenije}{A doubt}{Если у тебя есть хоть одно сомнение, то не спеши принимать решение.}{If you have even a single doubt, don’t be in a hurry to make the decision.}
\vocentry{1397}{Территория}{teri'torija}{A territory, an area}{Территория страны разделена на несколько провинций.}{The territory of the country is divided into several provinces.}
\vocentry{1398}{Тайна}{'tajna}{A secret; a mystery}{Это будет наша с тобой тайна. Никто не должен узнать об этом.}{It’s going to be our secret. Nobody must learn about it.}
\vocentry{1399}{Волк}{volk}{A wolf}{Ночью волк напал на стадо и убил несколько овец.}{A wolf attacked the herd at night and killed a few sheep.}
\vocentry{1400}{Терять}{te'rjat’}{To lose}{Я уезжаю из страны. Мне всё равно нечего терять.}{I’m leaving the country. I have nothing to lose anyway.}
\vocentry{1401}{Приятель}{pri'jatel’}{A fellow, a mate}{Он мой хороший приятель, но я не могу назвать его своим другом.}{He’s a good fellow of mine but I can’t call him my friend.}
\vocentry{1402}{Рукав}{ru'kav}{A sleeve}{Мне кажется, правый рукав этой рубашки короче, чем левый.}{It seems to me that the right sleeve of this shirt is shorter than the left one.}
\vocentry{1403}{Дедушка}{'dedushka}{A grandfather, a granddad}{Когда мы были детьми, дедушка много играл с нами.}{When we were kids, our grandfather played with us a lot.}
\vocentry{1404}{Вскочить}{vska'tchit’}{To leap (up, onto, to), to jump (up, onto), to get up quickly}{Не успел я вскочить на ноги, как соперник ударил меня снова.}{Hardly had I managed to leap to my feet when my opponent hit me again.}
\vocentry{1405}{Прочитать}{prachi'tat’}{To read, to read through (perfective)}{У меня не было времени прочитать книгу до конца. Можно мне оставить её ещё на день?}{I had no time to read the book till the end. May I keep it for another day?}
\vocentry{1406}{Честно}{'tchesnə}{Honestly, frankly; fair, fairly}{Честно говоря, я не в настроении идти на дискотеку.}{Honestly speaking, I’m not in the mood to go to the disco.}
\vocentry{1407}{Масло}{'maslə}{Butter; oil}{Закончилось масло. Добавь его в список покупок.}{We’ve run out of butter. Add it to the shopping list.}
\vocentry{1408}{Постоянный}{pasta'jannyj}{Constant, permanent}{Постоянный контроль родителей действует мне на нервы.}{My parents’ constant control is getting on my nerves.}
\vocentry{1409}{Всякая}{'fsjakaja}{All sorts of, all kinds of; any (feminine of ‘всякий)}{В этом ящике у меня хранится всякая всячина типа открыток, марок и старых билетов.}{In this drawer I keep all sorts of stuff like postcards, stamps and old tickets.}
\vocentry{1410}{Существование}{sustchestva'vanije}{Existence, being}{Я не верю в существование инопланетян. А ты?}{I don’t believe in the existence of aliens. And you?}
\vocentry{1411}{Лезть}{lest’}{To climb (up, onto); to intrude}{Пожарному пришлось лезть на дерево, чтобы снять кошку.}{The fireman had to climb up the tree to put the cat down.}
\vocentry{1412}{Младший}{'mladshij}{Younger; junior}{Мой младший брат ходит в детский сад.}{My younger brother goes to the kindergarten.}
\vocentry{1413}{Отправить}{at'pravit’}{To send (perfective)}{Скажите пожалуйста, какие бумаги мне нужно заполнить, чтобы отправить посылку за границу?}{Tell me please, what papers should I fill in to send a parcel abroad?}
\vocentry{1414}{Построить}{past'roit’}{To build, to construct (perfective)}{Мне интересно, как древние египтяне смогли построить пирамиды без современного оборудования.}{I wonder how ancient Egyptians managed to build the pyramids without modern equipment.}
\vocentry{1415}{Нервный}{'nervnyj}{Nervous}{Почему ты такой нервный сегодня? Что-нибудь случилось?}{Why are you so nervous today? Has anything happened?}
\vocentry{1416}{Пустить}{'pustit’}{To let, to allow, to permit; to let in (perfective)}{Я не могу пустить тебя на прогулку одну так поздно.}{I can’t let you go for a walk alone so late.}
\vocentry{1417}{Рождение}{razh'denije}{Birth}{Рождение нового человека – это удивительное событие в жизни каждой семьи.}{The birth of a new person is an amazing event in the life of every family.}
\vocentry{1418}{Заглянуть}{zaglja'nut’}{To peep in, to glance at; to drop by, to call on (colloquial)}{У моей бабушки был большой старый сундук. Мне всегда хотелось заглянуть в него.}{My grandmother had a large old chest. I always wanted to peep into it.}
\vocentry{1419}{Проговорить}{pragava'rit’}{To talk (for a long time)}{Мы могли проговорить всю ночь и не заметить этого.}{We could talk the whole night and not notice it.}
\vocentry{1420}{Воспоминание}{vaspəmi'nanije}{A memory, a recollection}{День нашей встречи – это моё любимое воспоминание.}{The day we met is my favorite memory.}
\vocentry{1421}{Проверить}{pra'verit’}{To check (perfective)}{Не забудь проверить, выключил ли ты свет перед уходом.}{Don’t forget to check if you’ve switched off the lights before leaving.}
\vocentry{1422}{Болеть}{ba'let’}{To be ill, to be sick; to cheer for}{Я ненавижу болеть летом, когда все плавают и загорают.}{I hate being ill in summer when everyone’s swimming and sunbathing.}
\vocentry{1423}{Тон}{ton}{A tone (of voice, music, color)}{Мне не нравится твой тон. Не разговаривай со мной так грубо.}{I don’t like your tone. Don’t talk to me so rudely.}
\vocentry{1424}{Понятие}{pan’'jatije}{A concept, an idea}{Попробуйте объяснить это научное понятие сами.}{Try to explain this scientific concept yourselves.}
\vocentry{1425}{Низкий}{'niskij}{Low; short (of a person)}{В этом городе очень низкий уровень преступности.}{This city has a very low crime rate.}
\vocentry{1426}{Далее}{'daleje}{Then; further, furthermore; farther}{Прочтите текст и далее ответьте на вопросы.}{Read the text and then answer the questions.}
\vocentry{1427}{Вечный}{'vechnyj}{Eternal, everlasting}{Я молюсь за вечный мир на земле.}{I pray for eternal peace on Earth.}
\vocentry{1428}{Песня}{'pjesnja}{A song}{О, это моя любимая песня! Сделай громче!}{Oh, that’s my favorite song! Turn it up!}
\vocentry{1429}{Морской}{mar'skoj}{Sea, seaside (adj)}{Дул свежий морской бриз, и мы гуляли по набережной.}{A fresh sea breeze was blowing and we were walking along the quay.}
\vocentry{1430}{Крыло}{kry'lo}{A wing}{Крыло птицы зажило, и мы выпустили её на волю.}{The bird’s wing healed and we let it out into the wild.}
\vocentry{1431}{Поднимать}{padni'mat’}{To lift, to raise, to elevate}{Тебе нельзя поднимать такие тяжёлые вещи. Дай я помогу тебе.}{You can’t lift such heavy things. Let me help you.}
\vocentry{1432}{Розовый}{'rozavyj}{Pink}{Розовый – цвет моды и гламура, но мне он не нравится.}{Pink is the color of fashion and glamour but I don’t like it.}
\vocentry{1433}{Нарушение}{naru'shenije}{Violation, breach}{Это нарушение прав человека, и с этим нельзя мириться.}{It’s the violation of human rights and it can’t be put up with.}
\vocentry{1434}{Осень}{'osen’}{Autumn}{Осень – удивительно красивая, красочная пора года.}{Autumn is an amazingly beautiful and colorful season of the year.}
\vocentry{1435}{Положенный}{pa'lozhennyj}{Due}{Постарайтесь подготовить все бумаги в положенный срок.}{Try to prepare all the papers in due time.}
\vocentry{1436}{Туман}{tu'man}{Fog}{Туман был таким густым, что мы ничего не видели.}{The fog was so dense that we couldn’t see anything.}
\vocentry{1437}{Ленинград}{lenin'grad}{Leningrad}{В советское время Санкт-Петербург назывался Ленинград.}{In Soviet times, St. Petersburg was called Leningrad.}
\vocentry{1438}{Платье}{'plat’je}{A dress}{Я купила новое платье, и теперь мне нужна подходящая сумочка.}{I’ve bought a new dress and now I need a handbag to match.}
\vocentry{1439}{Представление}{pretstav'lenije}{A performance}{Это было лучшее цирковое представление, которое я только видел!}{That was the best circus performance I’ve ever seen!}
\vocentry{1440}{Изменение}{izme'nenije}{A change, a modification, an alteration}{В правилах есть небольшое изменение. Будьте внимательны.}{There’s a slight change in the rules. Be attentive.}
\vocentry{1441}{Вариант}{vari'ant}{An option, a variant}{Посмотрите на этот диван. Это более доступный вариант.}{Have a look at this sofa. It’ a more affordable option.}
\vocentry{1442}{Шофёр}{sha'fjor}{A chauffeur, a driver}{Мой сосед такой богатый, что у него есть личный шофёр.}{My neighbor is so rich that he has a personal chauffeur.}
\vocentry{1443}{Опасный}{a'pasnyj}{Dangerous}{Снизь скорость – впереди опасный поворот.}{Slow down – there’s a dangerous turn ahead.}
\vocentry{1444}{Означать}{azna'tchat’}{To mean}{Я не могу понять, что бы могла означать её фраза. Она точно что-то скрывает от нас.}{I can’t understand what her phrase could mean. She’s surely hiding something from us.}
\vocentry{1445}{Вход}{fhot}{An entrance, an entry}{Убери эти коробки отсюда. Они загораживают вход в дом.}{Take these boxes away from here. They’re blocking the entrance to the house.}
\vocentry{1446}{Практически}{prak'titcheski}{Almost; virtually, practically}{Практически все наши сбережения были потрачены на восстановление дома после пожара.}{Almost all our savings were spent on the reconstruction of the house after the fire.}
\vocentry{1447}{Нынешний}{'nyneshnij}{Present, today’s}{Наш нынешний начальник гораздо более требовательный, чем предыдущий.}{Our present boss is much more demanding than the former one.}
\vocentry{1448}{Глупо}{'glupə}{Silly, foolishly}{Было глупо с твоей стороны обижаться на его шутку.}{That was silly of you to be offended by his joke.}
\vocentry{1449}{Национальный}{natsia'nal’nyj}{National}{На церемонии открытия спортсмены спели свой национальный гимн.}{At the opening ceremony, the sportsmen sang their national anthem.}
\vocentry{1450}{Чиновник}{tchi'novnik}{An official, a functionary}{Этот чиновник обвиняется во взяточничестве.}{This official is accused of bribery.}
\vocentry{1451}{Поверхность}{pa'verhnəst’}{A surface}{Поверхность этого кожаного чемодана гладкая и блестящая.}{The surface of this leather suitcase is smooth and glossy.}
\vocentry{1452}{Приятный}{pri'jatnyj}{Pleasant, pleasing, agreeable}{Приятный звук музыки помог мне успокоиться и расслабиться.}{The pleasant sound of the music helped me to calm down and relax.}
\vocentry{1453}{Каменный}{'kamennyj}{Stone (adj); stony}{Этот старинный каменный мост был построен в тринадцатом веке.}{This ancient stone bridge was built in the thirteenth century.}
\vocentry{1454}{Говориться}{gava'ritsa}{To say (impersonal), to be said, to be about}{В инструкции должно говориться, как правильно чистить прибор.}{The instructions must say how to clean the device correctly.}
\vocentry{1455}{Прочий}{'prochij}{Other}{В резюме следует указать ваше образование, предыдущее место работы и прочий опыт.}{You should state your education, previous employment and other experience in the CV.}
\vocentry{1456}{Испытывать}{is'pytyvat’}{To experience, to feel; to test}{Должно быть, сложно испытывать столько разных эмоций одновременно.}{It must be hard to experience so many different emotions at a time.}
\vocentry{1457}{Рано}{'ranə}{Early}{Я пришел слишком рано, и мне пришлось ждать, пока магазин откроется.}{I came too early and had to wait before the shop opens.}
\vocentry{1458}{Жениться}{zhe'nitsa}{To get married (of a couple when they marry each other) or to marry somebody (of a man when he takes a wife)}{Мы встречаемся ещё только несколько месяцев и не хотим жениться так скоро.}{We’ve been dating for a few months only and don’t want to get married so soon.}
\vocentry{1459}{Аппарат}{apa'rat}{An apparatus, a device; an apparat (political)}{Учёный изобрёл сложный аппарат и получил за него международную премию.}{The scientist invented a complex apparatus and received an international award for it.}
\vocentry{1460}{Церковь}{'tserkof’}{A church (both an institution and a building)}{Церковь всегда играла важную роль в истории России.}{The church has always played an important role in the history of Russia.}
\vocentry{1461}{Ручка}{'rutchka}{A handle, a knob; a pen}{Эта деревянная ручка нелепо смотрится на металлической двери.}{This wooden handle looks awkward on the metal door.}
\vocentry{1462}{Оглянуться}{aglja'nutsa}{To look back, to turn baсk (both literally and figuratively)}{Иногда приятно оглянуться и вспомнить прошлое.}{Sometimes it’s nice to look back and remember the past.}
\vocentry{1463}{Пригласить}{prigla'sit’}{To invite (perfective)}{Сколько гостей вы планируете пригласить на свадьбу?}{How many guests are you planning to invite to the wedding?}
\vocentry{1464}{Хозяйка}{ha'z’ajka}{A hostess; an owner (female); a wife (colloquial)}{Вечеринка была отличной, а хозяйка выглядела очаровательно.}{The party was great and the hostess looked charming.}
\vocentry{1465}{Строго}{'strogə}{Strictly}{Не смотри на меня так строго. Это всего лишь разбитая чашка.}{Don’t look at me so strictly. It’s just a broken cup.}
\vocentry{1466}{Принадлежать}{prinadle'zhat’}{To belong}{Этот дневник не может принадлежать твоему брату. На нём другое имя.}{This diary can’t belong to your brother. It has a different name on it.}
\vocentry{1467}{Внезапный}{vne'zapnyj}{Sudden, unexpected}{Внезапный ливень застал нас в поле, и нам негде было спрятаться.}{A sudden rain shower caught us in the field and we had nowhere to hide.}
\vocentry{1468}{Единый}{je'dinyj}{Unified, single; united}{Невозможно найти единый подход ко всем ученикам. Каждый из них – личность.}{It’s impossible to find a unified approach to all pupils. Each of them is a personality.}
\vocentry{1469}{Дворец}{dva'rets}{A palace}{Этот средневековый дворец – часть культурного наследия страны.}{This medieval palace is part of the cultural heritage of the country.}
\vocentry{1470}{Водитель}{va'ditel’}{A driver}{Водитель потерял управление и врезался в фонарный столб.}{The driver lost control of the car and crashed into a lamp post.}
\vocentry{1471}{Строить}{'stroit’}{To build, to construct}{Здесь собираются строить новый торговый центр.}{A new mall is going to be built here.}
\vocentry{1472}{Множество}{'mnozhestvə}{A multitude, a variety}{В мире множество животных и растений, о которых мы ничего не знаем.}{There’s a multitude of animals and plants in the world we don’t know anything about.}
\vocentry{1473}{Перейти}{perej'ti}{To cross, to get over; to pass on to}{Я не могу найти, где перейти дорогу. Где здесь пешеходный переход?}{I can’t find where to cross the road. Where’s a zebra crossing here?}
\vocentry{1474}{Столик}{'stolik}{A table (diminutive, a little table, very often in the context of booking a table in a restaurant)}{На день святого Валентина мы заказали столик в ресторане.}{On Valentine’s Day we booked a table in a restaurant.}
\vocentry{1475}{Урок}{u'rok}{A lesson}{Стандартный урок в русской школе длится 45 минут.}{A standard lesson in a Russian school lasts 45 minutes.}
\vocentry{1476}{Древний}{'drevnij}{Ancient}{Ежегодно этот древний греческий храм посещают тысячи туристов.}{Thousands of tourists visit this ancient Greek temple annually.}
\vocentry{1477}{Опасность}{a'pasnəst’}{Danger, risk}{Опасность курения очевидна, но, похоже, многих курильщиков это не волнует.}{The danger of smoking is evident but it looks like many smokers don’t care about it.}
\vocentry{1478}{Повезти}{paves'ti}{To be lucky, to have luck (impersonal)}{Как ей могло так повезти? Она опять выиграла в лотерею!}{How could she be so lucky? She’s won the lottery again!}
\vocentry{1479}{Французский}{fran'tsuskij}{The French language; French}{Французский считается языком любви.}{The French language is considered to be the language of love.}
\vocentry{1480}{Поглядеть}{paglja'det’}{To have a look at, to glance at (quickly, not for long)}{Можно поглядеть на вашего малыша?}{May I have a look at your baby?}
\vocentry{1481}{Горе}{'gore}{Misfortune, grief, trouble}{Её муж умер от сердечного приступа. Какое горе!}{Her husband died of a heart attack. What a misfortune!}
\vocentry{1482}{Честный}{'tchesnyj}{Honest; fair}{Я ценю твой честный ответ.}{I appreciate your honest answer.}
\vocentry{1483}{Способность}{spa'sobnast’}{An ability, a capability}{У тебя есть удивительная способность понимать меня в любой ситуации.}{You have an amazing ability to understand me in any situation.}
\vocentry{1484}{Повернуть}{paver'nut’}{To turn (both transitive and intransitive)}{На перекрёстке вам нужно будет повернуть направо, и вы сразу же увидите это здание.}{You’ll need to turn right at the crossroads and you’ll see the building at once.}
\vocentry{1485}{Остальное}{astal’'noje}{The rest}{Я ставлю свою семью на первое место, остальное для меня не важно.}{I place my family at the first place, the rest doesn’t matter to me.}
\vocentry{1486}{Добраться}{dab'ratsa}{To get to, to reach}{Мы не успеем добраться до аэропорта вовремя, если немного не поторопимся .}{We won’t manage to get to the airport in time if we don’t hurry a bit.}
\vocentry{1487}{Серьёзно}{ser’'joznə}{Seriously, earnestly}{Не воспринимай его слова так серьёзно. Он просто шутит.}{Don’t take his words so seriously. He’s just joking.}
\vocentry{1488}{Признать}{priz'nat’}{To admit; to acknowledge (perfective)}{Мне было нелегко признать свою вину, но я рад, что сделал это.}{It was uneasy for me to admit my guilt but I’m happy I did it.}
\vocentry{1489}{Поступить}{pastu'pit’}{To act, to behave; to enter, to join (perfective)}{Я не знаю, как поступить в этой ситуации. Ты можешь дать мне совет?}{I don’t know how to act in this situation. Can you give me advice?}
\vocentry{1490}{Работник}{ra'botnik}{A worker, an employee (blue-collar)}{Вы отличный работник, и, я думаю, вы заслуживаете повышение.}{You’re an excellent worker and I think you deserve a promotion.}
\vocentry{1491}{Долг}{dolg}{A duty; a debt}{Помочь родителям в старости – мой долг.}{Helping my parents in their old age is my duty.}
\vocentry{1492}{Специалист}{spetsia'list}{A specialist, an expert}{Нам нужен квалифицированный специалист с большим опытом.}{We need a qualified specialist with lots of experience.}
\vocentry{1493}{Торопиться}{tara'pitsa}{To hurry, to be in a hurry}{Не нужно торопиться – у нас много времени.}{There’s no need to hurry – we’ve got plenty of time.}
\vocentry{1494}{Итак}{i'tak}{So, then}{Итак, где мы остановились?}{So, where did we stop?}
\vocentry{1495}{Центральный}{tsen'tral’nyj}{Central}{Во время праздника центральный парк был переполнен людьми.}{During the fest the central park was crowded with people.}
\vocentry{1496}{Телевизор}{tele'vizər}{A TV set}{Давай купим современный телевизор с плоским экраном.}{Let’s buy a modern flat screen TV set.}
\vocentry{1497}{Зло}{zlo}{Evil, bad}{В сказках добро всегда побеждает зло.}{In fairy tales good always triumphs over evil.}
\vocentry{1498}{Вкус}{fkus}{A taste}{Я никогда не забуду вкус этих пирожных! Как ты их готовишь?}{I’ll never forget the taste of these cakes! How do you cook them?}
\vocentry{1499}{Муха}{'muha}{A fly (an insect)}{Что за надоедливая муха! Как нам её прихлопнуть?}{What a bothersome fly! How can we swat it?}
\vocentry{1500}{Поиск}{'poisk}{A search}{Простой поиск в Интернете докажет, что ты ошибаешься.}{A simple search on the Internet will prove that you’re wrong.}
\end{multicols}
